\scanpage{220}
\mspage{L 20}
% archive sheet number: 219
\deleted{par des \unsure{épithètes}} comme \char34{}vaseux\char34{}, \char34{}mou\char34{}, \char34{}inconsistant\char34{} - ou au bout opposé \inserted{\handwritten{comme}} \char34{}trivialités\char34{}, \char34{}enfantillage\char34{}, \char34{}bombinage\char34{} ...), mais où sera extirpé également ce travail même et ses fruits: celui où sont conçues, s'élaborent et naissent les notions et les visions nouvelles. Ce sera l'époque aussi où l'exercice de notre art sera réduit à d'arides et vaines voltiges de \char34{}poids et haltères\char34{} cérébraux, aux surenchères des prouesses pour \char34{}craquer\char34{} les problèmes au concours (\char34{}de difficulté proverbiale\char34{}) - l'époque d'une hypertrophie \handwritten{``}supermacho\handwritten{''} fiévreuse et stérile, prenant la suite de plus de trois siècles de renouvellement créateur.

\handwritten{7.} Mais à nouveau je digresse, en anticipant sur ce que la réflexion m'a enseigné. \deleted{\unsure{Je}}\inserted{\handwritten{J'étais}} part\deleted{\unsure{ais}}\handwritten{i} d'un double propos\handwritten{,} clairement présent en moi dès avant même les débuts de celle-ci: le propos d'une \char34{}déclaration d'intentions\char34{}, et (intimement lié à celui-ci, comme il vient d'apparaître) celui de m'exprimer au sujet de la nature du travail créateur. Il y avait pourtant un troisième propos encore, moins clairement présent sûrement au niveau conscient, mais répondant à un besoin p\deleted{\unsure{b}}lus profond et plus essentiel. Il était suscité par ces \char34{}interpellations\char34{} parfois déconcertantes, me parvenant de mon passé de mathématicien par la voix de ceux qui avaient été mes élèves ou mes amis((ou \deleted{les deux)}\inserted{du moins}, de \deleted{certains}\inserted{\handwritten{bon nombre}} d'entre eux). Au niveau épidermique, ce besoin se traduisait par une envie de \char34{}vider mon sac\char34{}, de dire quelques \char34{}vérités déplaisantes\char34{}. Mais plus profondément, sûrement, il y avait le besoin de \deleted{faire} \mbox{f\,a\,i\,r\,e} \mbox{c\,o\,n\,n\,a\,i\,s\,s\,a\,n\,c\,e} enfin avec un certain passé, que j'avais choisi jusque là d'éluder
C'est de ce besoin-là, avant tout, qu'est \deleted{\unsure{né}}\inserted{issu} Récoltes et Semailles. Cette longue réflexion a été ma \char34{}réponse\char34{}, au jour le jour, à cette pulsion de connaissance en moi, et à l'interpellation sans cesse renouvelée qui me venait du monde extérieur, du \char34{}monde mathématique\char34{} que j'avais quitté sans esprit de retour. M\unsure{à|a|a\deleted{i}}s à part les toutes premières pages de \char34{}Fatuité et Renouvellement\char34{} \deleted{\unsure{,}}\handwritten{(}celles qui en forment les deux premiers chapitres \char34{}Travail et découverte\char34{} et \char34{}Le rêve et le Rêveur\char34{}\deleted{\unsure{,}}\handwritten{)}, et dès le chapitre qui enchaîne \char34{}Naissance de la crainte\char34{} (p. 18), avec un \char34{}témoignage\char34{}qui n'était nullement prévu \handwritten{a}u programme, c'est ce besoin de faire connaissance de mon passé et de l'assumer pleinement, qui (je crois) a été la force principale en oeuvre dans l'écriture de Récoltes et Semail\handwritten{les.} \deleted{Récoltes et Semailles.}

L'interpellation qui m'était venue du monde des mathématiciens, et qui revenait sur moi avec une force nouvelle tout au cours \deleted{de l'écriture} de Récoltes et Semailles (et surtout, \inserted{\handwritten{au cours}} de l' \char34{}enquête\char34{} poursuivie dans les parties II et IV) avait pris d'emblée le masque de la suffisance ,
